\chapter*{前言：把注意力放回研究表达}
\addcontentsline{toc}{chapter}{前言}
\markboth{前言}{前言}

写论文时，我们经常在最后阶段才开始认真处理图片。数据已经分析，方法已经写完，结论似乎也很清楚，只差把几张图“画好看”。真正动手后，问题却接连出现：一张图塞不下全部信息，两条箭头容易被误解，几组数据放在一起掩盖了个体差异，屏幕上清晰的文字缩进论文后几乎读不出来。

这些困难与研究本身有关。作图要求我们重新决定：什么最重要，哪些内容属于证据，哪些属于解释，读者应该按什么顺序理解它们。图画不清楚，有时是布局还没有找到，有时则是要表达的问题还没有说清楚。

AI 让制作图片的起点变得更低。我们可以用文字描述一个场景，用草图指定构图，用已有图片解释想改哪里，也可以把数据交给可视化工具，尝试不同的表达方式。许多过去需要反复操作的步骤，如今可以通过一段说明开始。这给科研人员留下了更多比较方案和修正表达的空间。

与此同时，一张完成度很高的图片也更容易掩盖问题。错误的关系可以画得很流畅，遗漏的实验条件可以被整齐的版式遮住，没有依据的结论也可能配上一套看似专业的图形语言。我们需要学会把“能生成”进一步推进到“能解释、能核对、能交付”。

本书围绕这一过程展开。你将从一份绘图任务书开始，决定图的用途和核心问题，再选择数据图、方法图、机制图、实验图像或图形摘要的工作路线。我们会讨论构图、颜色、标签和图注，也会花时间处理修改、导出和源文件保存。每一部分都会说明 AI 可以承担哪些工作，以及哪些决定需要你根据研究材料来做。

全书使用 SciDraw AI 作为统一实践平台。统一平台能让读者持续处理同一个案例，减少在不同软件之间重新理解操作的负担。书中介绍的能力以实际界面和可完成的操作为依据；遇到需要下载、上传或人工补充的地方，会直接写出步骤。案例要求的功能尚未具备时，先给出能完成任务的路径，经过实现与实测的功能再进入教程。

你不需要先掌握一套绘画技巧。更有价值的准备，是带上一项熟悉的研究、一份保留完整的材料，以及愿意逐项检查图片的耐心。对科研绘图已经有经验的读者，可以直接从需要解决的问题进入：如何让图表保留研究设计，如何减少机制图中的歧义，如何把多个面板组织成一段可读的论证。

书中的练习会明确区分模拟材料与真实研究。模拟数据让我们能公开复现操作、讨论错误，并且不用把任何教学结果当成科学发现。真实案例则需要可靠来源和必要授权。希望你最终带走的，是一套能应用到自己研究中的工作方法：先确定表达，再使用工具，把每次重要的决定与其依据保留下来。

从下一章开始，请暂时把注意力从“生成一张漂亮图片”移到一个更具体的问题上：这张图要帮助读者理解什么？
